% \section{Summary of Findings}

% \section{Contributions to the Field}

% \section{Recommendations for Future Work}

The findings of our project on Synthetic Data Generation of Adverse Driving Conditions for Autonomous Driving Pretraining provide an important framework for overcoming the safety challenge in autonomous driving training for the diverse traffic conditions in
South Asia. Through our synthetic data creation and scalability efforts, we wish to successfully identify vehicle types like “Desi Nosimons” and auto-rickshaws and simulate dangerous driving scenarios based on areas which are not present in conventional Western datasets. Additionally, through our experiment, we were able to show that synthetic data was not merely a replacement for actual data collection efforts but were essential in order to cover all the important edge cases in order for the autonomous cars to be prepared for transitioning from test tracks to reality.
Looking towards the future, our focus would lie in evaluating our pre-trained models against real world benchmarks, namely through the use of BadODD for quantifying the reduction in the “Sim-to-Real” Gap. Lastly, our goal will be to improve the realism of our simulation to an extent that synthetic data will lead to the development of autonomous vehicles with enhanced safety and cultural compatibility in the Global South.
